%! TEX TS-program = xelatex
%----------------------------------------------------------------------------------------
%	Metropolia Thesis LaTeX Template
%----------------------------------------------------------------------------------------

%----------------------------------------------------------------------------------------
%	THESIS INFO
%----------------------------------------------------------------------------------------

% All general information (main language, title, author (you), degree programme, major
% option, etc.)
% Edit the file chapters/0info.tex to change these information

% Global information (title of your thesis, your name, degree programme, major, etc.)
\def\bilingual{yes}
\def\thesislang{finnish}
\def\secondlang{english}
\author{Arttu Pennanen}
%License
\def\thesiscopy{by}

%Finnish section, for title/abstract
\def\otsikko{Pragmaattisen funktionaalisen \\ohjelmoinnin arviointi}
\def\tutkinto{Insinööri (AMK)}
\def\kohjelma{Tieto- ja viestintätekniikka}
\def\suuntautumis{Ohjelmistotuotanto}
\def\thesisfi{Insinöörityö}
\def\ohjaajat{
    FM Simo Silander
}
\def\tiivistelma{
    TODO \newline

    Lorem ipsum dolor sit amet, consectetur adipiscing elit. Nunc interdum suscipit nunc. Sed suscipit vulputate nisi id vehicula. Proin mi sapien, tempor sit amet lacus vel, fringilla pulvinar purus. Duis non tempor elit. Donec lobortis est eget pretium faucibus. Maecenas sed convallis ligula. Nunc orci diam, ornare eget quam nec, tincidunt ultricies lacus. Curabitur in lorem ac elit facilisis malesuada id lacinia justo. Quisque placerat, metus quis suscipit interdum, nisl lectus auctor odio, ut consectetur eros leo ac est.
}
\def\avainsanat{funktionaalinen ohjelmointi, JavaScript, TypeScript, koodin luettavuus, koodin ylläpidettävyys, koodin suorituskyky, kehittäjäkokemus}
\def\aihe{TODO: Lyhyt kuvaus opinnäytetyöstä. Max 255 merkkiä. Funktionaalista koodia jne jne}%for the pdf metadata/properties. If not used, empty it and also the \def\subject.

%English section, for title/abstract
\title{Evaluating the Practical Application of Functional Programming}
\def\metropoliadegree{Bachelor of Engineering} % change to your needs, e.g. "master", etc.
\def\metropoliadegreeprogramme{Information and Communication Technology}
\def\metropoliaspecialisation{Software Engineering}
\def\thesisen{Bachelor’s Thesis} % change to your need, e.g. master's
\def\metropoliainstructors{
    Simo Silander, M. Sc.
}
\def\abstract{
    TODO: Abstract content. To force newline between paragraph in the abstract, you must have both a empty line and the newline command. \newline

    Lorem ipsum dolor sit amet, consectetur adipiscing elit. Nunc interdum suscipit nunc. Sed suscipit vulputate nisi id vehicula. Proin mi sapien, tempor sit amet lacus vel, fringilla pulvinar purus. Duis non tempor elit. Donec lobortis est eget pretium faucibus. Maecenas sed convallis ligula. Nunc orci diam, ornare eget quam nec, tincidunt ultricies lacus. Curabitur in lorem ac elit facilisis malesuada id lacinia justo. Quisque placerat, metus quis suscipit interdum, nisl lectus auctor odio, ut consectetur eros leo ac est.
}
\def\metropoliakeywords{functional programming, JavaScript, TypeScript, code readability, code maintainability, code performance, developer experience}
\def\subject{TODO: Short description of the thesis. Max 255 characters. Its the funksönal programming}%for the pdf metadata/properties. If not used, empty it and also the \def\aihe.


%----------------------------------------------------------------------------------------
%	GLOBAL STYLES
%----------------------------------------------------------------------------------------

% If you need extra package, etc. modify the style/style.tex file.
% If you are using Windows OS, you will need to change default font to Arial in that
% style/style.tex file (or install Liberation Sans font to your system).
% If you are using MacOS or linux, make sure you have Liberation Sans font installed.
% Global style. Normally should not be edited.
% If you use windows OS, eventually change \setmainfont to Arial
% Check around commit https://github.com/panunu/metropolia-thesis-latex/commit/a0c15ac77bab1a52c59c517a18080938e57bf5ef
% to see how the font files were manually added (after downloading them: https://pagure.io/liberation-fonts/ )

%\usepackage[l2tabu, orthodox]{nag}%check for obsolete packages (with outdated nag package?)

%condition e.g. for adding or not space in TOC,...
\usepackage{etoolbox}
\ifdefstring{\bilingual}{yes}{
	\usepackage[\secondlang,\thesislang]{babel}% finnish (or swedish) and english
}{
	\usepackage[\thesislang]{babel}% english only
}
\usepackage{iflang}
\usepackage{amsmath}
\usepackage{amsfonts}%extra mathematical symbols
\usepackage{amssymb}
\usepackage{fontspec}
\usepackage{titlesec}
\usepackage{mathtools}%enhance the appearance of documents containing a lot of mathematics
\usepackage{siunitx}  % Load siunitx package for SI units
\usepackage[version=3]{mhchem}
\usepackage{tikz} % mindmaps, flowcharts, piecharts, examples at http://www.texample.net/tikz/examples/
\usetikzlibrary{shapes.geometric, arrows}

\usepackage{ragged2e}
\IfLanguageName{finnish}{
	\RaggedRight%2021 template, align left and hyphennization for Finnish version
	\setlength{\RaggedRightRightskip}{0pt plus 1fil}%TODO: fix the Overfull/Underfull warnings when \RaggedRight (likely in title and abstact)
}{}
\makeatletter
\let\@gnewline\@raggedtwoe@saved@gnewline% Restore original functionality of \newline
\let\\\@raggedtwoe@savedcr% Restore original functionality of \\
\makeatother

%celsius
\newcommand{\celsius}{\ensuremath{^\circ\mathrm{C}}}
%for compact list
\usepackage{enumitem}
%\usepackage{verbatim}%for block comment
%forcing single line spacing in bibliography
\DisemulatePackage{setspace}
\usepackage{setspace}
%\usepackage{eurosym}%euro symbol
%try to count
\usepackage{totcount}
%insert source code
%require -8bit -shell-escape in the xelatex compile command
%if compiling locally, consider options cachedir=minted,outputdir=~/.tex
\usepackage[newfloat]{minted}
\setminted{tabsize=2,linenos,breaklines,breaksymbolleft={\quad},baselinestretch=1}
\setmintedinline{breaklines,breakanywhere}
\usepackage[singlelinecheck=false]{caption}
%force the width of a table instead of column
\usepackage{tabularx}
\usepackage{booktabs} %why not booktabs? :3

\usepackage{csquotes}% avoid warning with babel

\usepackage{float} % For forced figure location with modifier H (\begin{figure}[H])
\usepackage{wrapfig}

%set date format to D.M.YYYY
\def\pvm{\the\day.\the\month.\the\year}
%set date format to D Month YYYY
\usepackage[en,useregional=false]{datetime2}
\DTMsetup{datesep=.}
\DTMnewdatestyle{dMonthyyyy}
{%
	\renewcommand{\DTMdisplaydate}[4]{%
		\number##3 % day
		~% separator
		\DTMenglishmonthname{##2}% month name
		~% separator
		\number##1% year
	}%
	\renewcommand{\DTMDisplaydate}[4]{%
		\DTMdisplaydate{##4}{##3}{##2}{##1}%
	}%
}
\DTMsetdatestyle{dMonthyyyy}
\date{\today}


\usepackage{xcolor}
\usepackage{tcolorbox}
\tcbuselibrary{listingsutf8} % For including monospace font settings

\newcommand{\key}[1]{%
	\begin{tcolorbox}[
			colback=gray!20, % Light gray background
			colframe=gray!20, % Same as background, effectively invisible
			boxrule=0pt, % No border
			arc=4pt, % Rounded corners
			left=2pt, % Left padding
			right=2pt, % Right padding
			top=2pt, % Top padding
			bottom=2pt, % Bottom padding
			boxsep=0pt, % No separation between text and frame, padding only
			height=12pt,
			valign=center,
			nobeforeafter, % No additional space before and after the box
			tcbox raise base, % Aligns the box with the baseline of the surrounding text
			hbox % Ensures the box is only as wide as its contents
		]
		\texttt{\small #1} % Smaller monospaced font for the key text
	\end{tcolorbox}
}

\newcommand\tn[1]{\textnormal{#1}} %use \tn instead of \textnormal
\newcommand\reaction[1]{\begin{equation}\ce{#1}\end{equation}} %\reaction{} for chemical reactions

%NORMAL TEXT
%all text, title, etc. in the same font: Arial
%NOTE: fontname is case-sensitive
\setmainfont[Scale=0.98]{Liberation Sans}
%line space
\linespread{1.46}
\AtBeginEnvironment{tabular}{\singlespacing}
%\doublespacing
%margin
%geometry moved after hyperref
\setlength{\parindent}{0pt} %first line of paragraph not indented
\setlength{\parskip}{16.5pt} %one empty line to separate paragraph
%list with small line space separation
\tightlists
\setlist[itemize]{before=\singlespacing,itemsep=6pt,leftmargin=63pt,labelsep=22pt,topsep=1.5pt,partopsep=0pt,after=\vspace{-22pt}\newline}
\setlist[enumerate]{before=\singlespacing,itemsep=6pt,leftmargin=63pt,labelsep=22pt,topsep=1.5pt,partopsep=0pt,after=\vspace{-22pt}\newline}

%IMAGE - FIGURE
%the figures should be placed in the "figures" folder
\graphicspath{{figures/}}
%figure number without chapter (1.1, 1.2, 2.1) to (1, 2, 3)
\counterwithout{figure}{chapter}
%border around images
\setlength\fboxsep{0pt}
\setlength\fboxrule{0.5pt}
%space after figure caption (and other float elements)
\setlength{\belowcaptionskip}{-7pt}
\setlength{\intextsep}{16.5pt}%space around floats
\AtBeginEnvironment{table}{\addvspace{16.5pt}}

%TABLE
\counterwithout{table}{chapter}

%SOURCE CODE
\newenvironment{code}
{
	\begin{samepage} % Start keeping content on the same page
		\captionsetup{type=listing}
		}
		{
	\end{samepage} % End keeping content on the same page
}

\IfLanguageName{finnish}
{
	\SetupFloatingEnvironment{listing}{name=Koodiesimerkki}
}
{}%was Listaus
%\counterwithout{lstlisting}{chapter}
%moved after begin document, otherwise does not compile

%TOC
%change toc title
\IfLanguageName{finnish}{\addto{\captionsfinnish}{\renewcommand*{\contentsname}{Sisällys}}}{}
%remove dots
\renewcommand*{\cftdotsep}{\cftnodots}
%chapter title and page number not in bold
\renewcommand{\cftchapterfont}{\normalfont}
\renewcommand{\cftchapterpagefont}{\normalfont}
%sub section in toc
\setcounter{tocdepth}{2}
%subsection numbered
\setcounter{secnumdepth}{2}
\renewcommand{\tocheadstart}{}
\renewcommand{\printtoctitle}[1]{\fontsize{13.5pt}{13.5pt}\bfseries #1\vspace*{-4pt}}
%\renewcommand{\aftertoctitle}{\vspace*{-22pt}\afterchaptertitle}
%spacing after a chapter in toc
\preto\section{%
  \ifnum\value{section}=0\addtocontents{toc}{\vskip11pt}\fi
 }
%spacing after a section in toc
\renewcommand{\cftsectionaftersnumb}{\vspace*{-3pt}}
%spacing after a subsection in toc
\renewcommand{\cftsubsectionaftersnumb}{\vspace*{-1pt}}
%appendix in toc with "Appendix " + num
\IfLanguageName{finnish}{
	\renewcommand*{\cftappendixname}{Liite\space}
	\renewcommand{\appendixtocname}{Liitteet}
}{\renewcommand*{\cftappendixname}{Appendix\space}}
%appendix header
\IfLanguageName{finnish}{\def\appname{Liite\space}}{\def\appname{Appendix\space}}

%TITLES
%chapter title
%\clearforchapter{\clearpage}
\titleformat{\chapter}
{\fontsize{14pt}{14pt}\bfseries\linespread{1}}%\clearpage
{\thechapter}{.5cm}{}
\titlespacing*{\chapter}{0pt}{-.42cm}{.5pt}
\titleformat{\section}
{\fontsize{13.5pt}{13.5pt}\normalfont}
{\thesection}{.5cm}{}
\titlespacing*{\section}{0pt}{9pt}{0pt}
\titleformat{\subsection}
{\fontsize{12.7pt}{12.7pt}\normalfont}
{\thesubsection}{.5cm}{}
\titlespacing*{\subsection}{0pt}{11pt}{0pt}


%QUOTE
\renewenvironment{quote}
{\list{}{\rightmargin=0pt\leftmargin=2.2cm\topsep=-14pt}%
	\item\relax\singlespacing}%\fontsize{10pt}{10pt}
{\vspace{8pt}\endlist}

%BIBLIOGRAPHY
%TODO: toggle for Harward v.s. Vancouver
\usepackage[
	backend=biber,
	bibencoding=utf8,%
	citetracker=true,%
	isbn=false,%
	doi=false,%
	url=true,%
	usetranslator=true,%
	bibstyle=ext-authoryear,%
	citestyle=numeric-comp,%
	sorting=none,%
	sortcites=true,%
	block=none,%
	terseinits=false,%
	giveninits=false,%
	maxbibnames=99,%
]{biblatex}

\setlength{\bibitemsep}{11pt}
\setlength{\biblabelsep}{1cm}%with 1cm horizontal gap

\makeatletter
\RequireBibliographyStyle{numeric}
\makeatother

% set right format
\renewcommand*{\multicitedelim}{;\space}
\renewcommand*{\multinamedelim}{;\space}
\renewcommand*{\finalnamedelim}{~\&\space}
\DeclareFieldFormat{biblabeldate}{#1}
\DeclareDelimFormat[bib]{nameyeardelim}{\addperiod\space}
%\DeclareNameAlias{sortname}{last-first} % deprecated
\DeclareNameAlias{default}{family-given}
\DeclareFieldFormat{labelnumberwidth}{#1} % remove () from label number
\DeclareFieldFormat{title}{#1} % normal font for title
\DeclareFieldFormat{citetitle}{#1}
\DeclareFieldFormat{journaltitle}{#1} % remove underline
\DeclareFieldFormat*{url}{%
	\ifentrytype{online}{\IfLanguageName{finnish}{Verkkoaineisto}{Online}\addperiod\space}{}%
	\ifentrytype{video}{Video\addperiod\space}{}%
	\ifentrytype{podcast}{Podcast\addperiod\space}{}%
	\textless\url{#1}\textgreater\addperiod
} % you can modify how to url looks here

\DeclareFieldFormat{urldate}{%
	\ifentrytype{video}{\space Katsottu\space#1}{%
		\ifentrytype{podcast}{\space Kuunneltu\space#1}{\space Luettu\space#1}}%
}

\DefineBibliographyStrings{finnish}{%
	urlfrom = {},%
	fromjapanese = {japanista},%
	fromenglish = {englannista},%
	fromfinnish = {suomesta},%
	fromgerman = {saksasta},%
}
{
	%new cite command: "Vancouver Short"
	\DeclareCiteCommand{\citeVS}
	{\usebibmacro{prenote}}
	{\usebibmacro{author}, \usebibmacro{title}}
	{\multicitedelim}
	{\usebibmacro{postnote}}

	% new cite command: "Vancouver Short Collection" - necessary when referencing whole collections.
	\DeclareCiteCommand{\citeVSc}
	{\usebibmacro{prenote}}
	{\usebibmacro{editor}, \usebibmacro{title}}
	{\multicitedelim}
	{\usebibmacro{postnote}}
}

\newcommand{\engl}[1]{engl. \textit{#1}.}

\DeclareCiteCommand{\citep}[\mkbibbrackets]
{\usebibmacro{cite:init}%
	\usebibmacro{prenote}}
{\usebibmacro{citeindex}%
	\usebibmacro{cite:comp}}
{}
{\usebibmacro{cite:dump}%
	\usebibmacro{postnote}\addperiod}


\addbibresource{biblio.bib}% for biblatex you need out \printbibliography too

%count the appendices (since the chapter counter is reset after \appendix).
%! require to complie 2 times
\regtotcounter{chapter}

% metadata (title, author, lang,...) for accessibility, etc.
\usepackage{hyperref}
\usepackage{hyperxmp}
\def\isolang{\IfLanguageName{finnish}{fi}{en}} %iso code (based on main language)
\ifdefstring{\thesiscopy}{all}{%
	\def\copyen{Copyright \textcopyright\ \the\year{}, \theauthor}
	\def\copyfi{Kaikki oikeudet pidätetään.  \textcopyright\ \the\year{}, \theauthor}
}{%
	\usepackage[type={CC},modifier={\thesiscopy},version={4.0},]{doclicense}
	\def\copyen{\thetitle\ \textcopyright\ \the\year{} by \theauthor\ is licensed under {\doclicenseLongNameRef}}
	\def\copyfi{\otsikko\ \textcopyright\ \the\year{}, jonka tekijä on \theauthor, on julkaistu {\doclicenseLongNameRef} lisenssillä}
}
\hypersetup{%
	pdfdisplaydoctitle,
	%breaklinks=true,%searching for overfull warnings
	pdfencoding=auto,
	bookmarksdepth=subsection,
	unicode=true,
	keeppdfinfo=true,
	pdflang={\isolang},
	pdfmetalang={\isolang},
	pdftitle={\IfLanguageName{finnish}{\otsikko}{\thetitle}},
	pdfkeywords={\IfLanguageName{finnish}{\avainsanat}{\metropoliakeywords}},
	pdfcopyright={\IfLanguageName{finnish}{\copyfi}{\copyen}},
	pdfsubject={\IfLanguageName{finnish}{\aihe}{\subject}},
}

\ifdefstring{\bilingual}{yes}{%metadata (title and copyright) in multiple language
	\XMPLangAlt{\IfLanguageName{finnish}{en}{fi}}{%
		pdftitle={\IfLanguageName{finnish}{\thetitle}{\otsikko}},
		pdfcopyright={\IfLanguageName{finnish}{\copyen}{\copyfi}},
		pdfsubject={\IfLanguageName{finnish}{\subject}{\aihe}},
	}
}{}
\urlstyle{same}

%moved after hyperred as can cause conflicts
\usepackage{pdfcomment}%try the alt text for graphics
\usepackage{accsupp}
\newcommand{\AltText}[2]{\BeginAccSupp{method=pdfstringdef,unicode,Alt={{#1}}}\pdftooltip{{#2}}{{#1}}\EndAccSupp{}}

\usepackage[top=2.5cm, bottom=3cm, left=4cm, right=2cm, nofoot]{geometry}

\usepackage{pgfplots} %simple plots etc
\pgfplotsset{compat=1.16}
\usepackage{pgfplotstable}

% Abbreviations, acronym and glossary, in case of bug, remove temporary the noredefwarn
\usepackage[acronym,toc,nonumberlist,section=chapter,noredefwarn]{glossaries}%xindy,%toc, ,nomain
\newglossarystyle{mystyle}{%
	\setglossarystyle{list}% base this style on the list style
	\renewcommand*{\glossentry}[2]{%
		\item[\glsentryitem{##1}%
		            \glstarget{##1}{\glossentryname{##1}:}]
		\glossentrydesc{##1}\glspostdescription\space ##2}
}
\setglossarystyle{mystyle}

\renewcommand*{\glsclearpage}{}

% Normally, you do not need to modify the title style. It's content comes from the
% chapters/0info.tex file.
% TITLE PAGE
% Normally, you should not edit this file.

\makeatletter
\renewcommand{\maketitle}{
\newgeometry{left=4cm,top=2.7cm}
\thispagestyle{empty}

\includegraphics[width=9.1cm,keepaspectratio]{Metropolia_L_oranssi}

\IfLanguageName{finnish}{
  \vspace*{1.36cm}
}{
  \vspace*{1.18cm}
}
\tn{\large\theauthor\\[9.2mm]{\fontsize{27.5pt}{27.5pt}\normalfont\color[RGB]{155,50,35}\IfLanguageName{finnish}{\otsikko}{\thetitle}}}%\\[22pt]\LARGE\alaotsikko\\[1.75cm]}

\vspace*{2.38cm}
{\setstretch{1.512}
  \IfLanguageName{finnish}{
    Metropolia Ammattikorkeakoulu\\
    \tutkinto \\
    \kohjelma \\
    \thesisfi\\
    \pvm % D.M.YYYY date format for Finnish.
  } {
    Metropolia University of Applied Sciences\\
    \metropoliadegree \\
    \metropoliadegreeprogramme \\
    \thesisen\\
    \@date %D Month YYYY for English
  }
}
\clearpage
\restoregeometry
}
\makeatother



%----------------------------------------------------------------------------------------
%	ABBREVIATION AND GLOSSARY
%----------------------------------------------------------------------------------------

% Add/edit all your acronyms, abbreviations, glossary entries, etc. definitions in
% chapters/0abbr.tex file.
% You can have as many as you wish. Only the ones you use in your text (inserted with
% \gls{} command) will print in the Glossary/Lyhenteet.
% Generate the glossary
\makeglossaries

% Acronyms, abbreviations, etc.


% \newacronym{fp}{fp}{Functionaalinen ohjelmointi}

% Glossary entries


\newglossaryentry{ts}{
	name={TypeScript},
	description={ohjelmointikieli (ts)}
}

\newglossaryentry{js}{
	name={JavaScript},
	description={ohjelmointikieli (js)}
}

\newglossaryentry{functional_programming}{
	name={Funktionaalinen ohjelmointi},
	description={Ohjelmointiparadigma, joka korostaa laskennan mallintamista funktioiden avulla}
}

\newglossaryentry{programming_paradigm}{
	name={Ohjelmointiparadigma},
	description={Ohjelmoinnin lähestymistapa tai malli, joka määrittelee, miten ohjelmat rakennetaan ja miten ohjelmointiongelmia ratkaistaan. Esimerkkejä tunnetuista ohjelmointiparadigmoista ovat funktionaalinen, olio- ja imperatiivinen ohjelmointi}
}

\newglossaryentry{imperative_programming}{
	name={Imperatiivinen ohjelmointi},
	description={Ohjelmointiparadigma, jossa ohjelmoija antaa tarkat ohjeet siitä, kuinka tehtävä suoritetaan vaihe vaiheelta. Tämä lähestymistapa keskittyy ohjelman tilan hallintaan ja muuttamiseen käskyjen avulla, kuten muuttujien asettaminen ja silmukoiden käyttö}
}

\newglossaryentry{declarative_programming}{
	name={Deklaratiivinen ohjelmointi},
	description={Ohjelmointiparadigma, jossa ohjelmoija määrittelee, mitä lopputuloksen tulisi olla, mutta ei yksityiskohtaisesti sitä, miten tämä tulos saavutetaan}
}

\newglossaryentry{object_oriented_programming}{
	name={Olio-ohjelmointi},
	description={Ohjelmointiparadigma, joka perustuu olioiden, eli ohjelmakomponenttien, ympärille, jotka yhdistävät sekä dataa että siihen liittyviä toimintoja. Olio-ohjelmointi korostaa perinnän, kapseloinnin ja polymorfismin kaltaisia periaatteita}
}

\newglossaryentry{set_theory}{
	name={Joukko-oppi},
	description={\protect\engl{set theory} Joukko-oppi on joukkojen ominaisuuksiin perehtynyt matematiikan osa-alue. Joukko-oppi on perustavanlaatuisessa merkityksessä tietorakenteissa. Joukko-opissa olennaisimpia laskutoimituksia ovat esimerksiksi unioni, leikkaus ja joukkoerotus}
}

\newglossaryentry{category_theory}{
	name={Kategoriateoria},
	text={kategoriateoria},
	description={\protect\engl{category theory} Kategoriateoria on matematiikan osa-alue, joka tutkii erittäin yleistävällä tasolla rakenteiden välisiä suhteita ja yhteyksiä. Kategoriateoriaa käytetään funktionaalisen ohjelmoinnin apukeinona}
}

\newglossaryentry{monad}{
	name={Monadi},
	plural={Monadit},
	description={\protect\engl{monad} Monadi on abstrakti tietotyyppi, jota käytetään funktionaalisessa ohjelmoinnissa.}
}

\newglossaryentry{language_agnostic}{
	name={Kieliagnostinen},
	description={\protect\engl{language agnostic} Kieliagnostisuudella tarkoitetaan lähestymistapaa, jossa ei keskitytä tiettyyn ohjelmointikieleen, vaan tarkastellaan peruskäsitteitä ja -periaatteita yleisellä tasolla, sovellettavissa eri kieliin riippumatta niiden syntaksista tai erityispiirteistä}
}

\newglossaryentry{correctness}{
	name={Koodin oikeellisuus},
	description={\protect\engl{correctness} Koodin oikeellisuus on ohjelman kyky täyttää sille asetetut vaatimukset virheettömästi. Toisaalta ilman vaatimuksia koodi ei voi olla oikeellista}
}

\newglossaryentry{pure_function}{
	name={Puhdas funktio},
	description={\protect\engl{pure function} Funktiot, jotka eivät vaikuta ohjelman tilaan ja palauttavat aina saman tuloksen samoilla syötteillä}
}

\newglossaryentry{immutable_data}{
	name={Muuttumaton data},
	description={\protect\engl{immutable data} Data, jota ei voi muuttaa sen luomisen jälkeen}
}

\newglossaryentry{higher_order_function}{
	name={Korkeamman asteen funktio},
	plural={Korkeamman asteen funktiot},
	description={\protect\engl{higher order function} Funktiot, jotka voivat ottaa toisia funktioita syötteenään tai palauttavat funktioita tuloksena}
}

\newglossaryentry{combinator}{
	name={Kombinaattori},
	description={\protect\engl{combinator} Funktionaalisen ohjelmoinnin rakennuspalikka, joka yhdistää funktioita ilman muuttujien käyttöä},
	plural={Kombinaattorit}
}

\newglossaryentry{composed_function}{
	name={Yhdistetty funktio},
	plural={Yhdistetyt funktiot},
	description={\protect\engl{function composition} Useammasta funktiosta koostettu funktio, jossa funktiot suoritetaan toinen toisensa jälkeen niin, että edellisen tuloste anneetaan seuraavan syötteeksi}
}

%----------------------------------------------------------------------------------------
%	DOCUMENT STARTS HERE...
%----------------------------------------------------------------------------------------

\begin{document}
\IfLanguageName{finnish}{
}{
  \raggedright%2021 template, align left, no hyphennization for English version
}
\counterwithout{listing}{chapter}

%----------------------------------------------------------------------------------------
%	TITLE PAGE
%----------------------------------------------------------------------------------------

\input{style/title_headers.tex}
\maketitle
\newpage

%----------------------------------------------------------------------------------------
%	ABSTRACT / Tiivistelmä
%----------------------------------------------------------------------------------------

% If you are international student writing in English, ignore the Finnish abstract.
% If you are Finnish citizen, you must have 2 abstracts, one in Finnish (or Swedish
% depending on your mother tongue) and one in English regardless of the main language of
% your thesis. Normally, you do not need to modify the abstract style. It's content comes
% from the chapters/0info.tex file.
\ifdefstring{\bilingual}{no}{%
    \input{style/abstract_en.tex}
    }{%
    \IfLanguageName{finnish}{%order of abstracts based on main language and spacing hell
        % Abstract in Finnish
% Normally, you should not edit this file. Everything comes from chapter/0info.tex

\pagestyle{empty} %remove page number
\newgeometry{top=1.5cm,left=4cm,right=.5cm}
\begin{otherlanguage}{finnish}
{\large\textbf{Tiivistelmä}}
{\renewcommand{\arraystretch}{1.1}
  \begin{tabular}{@{}p{4.7cm} >{\raggedright\arraybackslash}p{10.8cm}@{}}
  Tekijä: & \makeatletter\@author\makeatother
  \\
  Otsikko: & \otsikko
  \\
  Sivumäärä: & \pageref*{LastPage} sivua %+ \total{chapter} liitettä
  \\
  Aika: & \pvm
  \\[6.5mm]
  Tutkinto: & \tutkinto
  \\
  Tutkinto-ohjelma: & \kohjelma
  \\
  Ammatillinen pääaine: & \suuntautumis
  \\
  Ohjaajat: & \ohjaajat
  \\[11mm]
  \cmidrule[.7pt](l{-.15em}r{5.5em}){1-2}
  \multicolumn{2}{>{\raggedright}p{15.5cm}}{
  \vspace{.2mm}
  \tiivistelma
  }
  \\
  \\
  Avainsanat: & \avainsanat
  \\
\end{tabular}
}
\end{otherlanguage}
\restoregeometry
\clearpage


        % Abstract in English
% Normally, you should not edit this file. Everything comes from chapter/0info.tex

\pagestyle{empty} %remove page number
\newgeometry{top=1.77cm,left=4cm,right=.5cm}
\begin{otherlanguage}{english}
  {\large\textbf{Abstract}}
  {\renewcommand{\arraystretch}{1.1}
    \begin{tabular}{@{}p{4.7cm} >{\raggedright\arraybackslash}p{10.8cm}@{}}
      Author:             & \makeatletter\@author\makeatother
      \\
      Title:              & \makeatletter\@title\makeatother
      \\
      Number of Pages:    & \pageref*{LastPage} pages         %+ \total{chapter} appendices
      \\
      Date:               & \makeatletter\thedate\makeatother
      \\[6.5mm]
      Degree:             & \metropoliadegree
      \\
      Degree Programme:   & \metropoliadegreeprogramme
      \\
      Professional Major: & \metropoliaspecialisation
      \\
      Supervisors:        & \metropoliainstructors
      \\[9mm]
      \cmidrule[.7pt](l{-.15em}r{5.5em}){1-2}
      \\
      \multicolumn{2}{>{\raggedright}p{15.5cm}}{
      \makebox[0pt][l]{\hspace{-0.2cm}\parbox{\dimexpr\textwidth-0.5cm\relax}{\abstract}}
      }
      \\
      \\
      Keywords:           & \metropoliakeywords
      \\
    \end{tabular}
  }
\end{otherlanguage}
\restoregeometry
\clearpage


        }{
        \input{style/abstract_en.tex}
        % Abstract in Finnish
% Normally, you should not edit this file. Everything comes from chapter/0info.tex

\pagestyle{empty} %remove page number
\newgeometry{top=2.04cm,left=4cm,right=.5cm}
\begin{otherlanguage}{finnish}
{\large\textbf{Tiivistelmä}}
{\renewcommand{\arraystretch}{1.1}
  \begin{tabular}{@{}p{4.7cm} >{\raggedright\arraybackslash}p{10.8cm}@{}}
  Tekijä: & \makeatletter\@author\makeatother
  \\
  Otsikko: & \otsikko
  \\
  Sivumäärä: & \pageref*{LastPage} sivua %+ \total{chapter} liitettä
  \\
  Aika: & \pvm
  \\[6.5mm]
  Tutkinto: & \tutkinto
  \\
  Tutkinto-ohjelma: & \kohjelma
  \\
  Ammatillinen pääaine: & \suuntautumis
  \\
  Ohjaajat: & \ohjaajat
  \\[4mm]
  \cmidrule[.7pt](l{-.15em}r{5.5em}){1-2}
  \\
  \multicolumn{2}{>{\raggedright}p{15.5cm}}{
  \tiivistelma
  }
  \\
  \\
  Avainsanat: & \avainsanat
  \\
\end{tabular}
}
\end{otherlanguage}
\restoregeometry
\clearpage


    }
}
%----------------------------------------------------------------------------------------
%	License? Acknowledgement?
%----------------------------------------------------------------------------------------

% Uncomment next line and edit chapters/0license.tex if you want license in your thesis.
% License of your thesis
% If you wish to explain what it means. When you publish your thesis in https://theseus.fi
% you will be able to choose between some Creative Commons licenses
% https://creativecommons.org
% Adapt this example text to your taste.
% This would also be the right place to explain the license you choose for the code you
% produced for your thesis.

\pagestyle{empty}
\newgeometry{top=2.04cm,left=4cm}

{\large\textbf{Lisenssit}}

\ifdefstring{\thesiscopy}{all}{}{%
  \begin{wrapfigure}{r}{0.3\textwidth}
    \doclicenseImage
  \end{wrapfigure}
}
\IfLanguageName{finnish}{\copyfi}{\copyen}


\textbf{Voit vapaasti:}
\begin{itemize}
  \item Jakaa {\textemdash} kopioida aineistoa ja levittää sitä edelleen missä tahansa välineessä ja muodossa missä tahansa tarkoituksessa, myös kaupallisesti.
  \item Muunnella {\textemdash} remiksata ja muokata aineistoa sekä luoda sen pohjalta uusia aineistoja missä tahansa tarkoituksessa, myös kaupallisesti.
\end{itemize}

\textbf{Seuraavilla ehdoilla:}
\begin{itemize}
  \item Nimeä {\textemdash} Sinun on mainittava lähde asianmukaisesti, tarjottava linkki lisenssiin sekä merkittävä, mikäli olet tehnyt muutoksia.\\ Voit tehdä yllä olevan millä tahansa kohtuullisella tavalla, mutta et siten, että annat ymmärtää lisenssinantajan suosittelevan sinua tai teoksen käyttöäsi.
  \item Ei muita rajoituksia {\textemdash} Et voi asettaa sellaisia oikeudellisia ehtoja tai teknisiä estoja, jotka estävät oikeudellisesti muita tekemästä mitään sellaista, minkä lisenssi sallii.
\end{itemize}
\bigskip
{\large\textbf{PDF-tiedoston linkit}}

Navigoinnin helpottamiseksi PDF-tiedostossa on käytössä linkkejä. PDF-katseluohjelmasta riippuen linkeissä näkyy ääriviivat:

\begin{itemize}
  \item\fcolorbox{red}{white}{Punaiset ääriviivat}: tiedoston sisäiset linkit sanastoon
  \item\fcolorbox{green}{white}{Vihreät ääriviivat}: tiedoston sisäiset linkit lähdeluetteloon
  \item \fcolorbox{cyan}{white}{Siniset ääriviivat}: ulkoiset linkit verkkosivuille.
\end{itemize}

\noindent Sisäisistä linkeistä voi palata takaisin näppäinyhdistelmällä \key{Alt}+ \key{Vasen nuoli}
(tai \key{Cmd}+ \key{Vasen nuoli}).
\restoregeometry
\clearpage


% Uncomment next line and edit chapters/0acknowledgement.tex if you want acknowledgements.
%\input{chapters/0acknowledgement.tex}

%----------------------------------------------------------------------------------------
%	TABLE OF CONTENTS
%----------------------------------------------------------------------------------------

% Table of Content
% Normally, you don't have to modify this file.
% In case spacing of appendix/liite goes crazy, look in ../thesis.tex and look
% for the hack described in APPENDICES.

\makeevenhead{plain}{}{}{}
\makeoddhead{plain}{}{}{}
\pagestyle{empty} %remove page number in toc (if longer than 2 pages)
\newgeometry{top=2.04cm, left=4cm}
\tableofcontents*
\pagestyle{empty} %remove page number in toc (if longer than 1 pages)

\restoregeometry
\clearpage



%list of figure, tables would come here if relevant?

%----------------------------------------------------------------------------------------
%	Lyhenteet / Abbreviation
%----------------------------------------------------------------------------------------

% If you don't use abbreviations/glossary, remove the following line.
% Abbreviation and Glossary
% Normally, you don't have to modify this file. Your abbreviations, etc. goes in
% ../chapters/0abbr.tex file.

\begin{singlespacing}

	% \gsladdall%would add all terms even if not used in your text.
	\addtocontents{toc}{\cftpagenumbersoff{chapter}}
	{
		\titleformat{\section}
		{\fontsize{13.5pt}{13.5pt}\bfseries\normalfont}
		{\thesection}{.5cm}{}
		\renewcommand*{\glossarypreamble}{\vspace{.7\baselineskip}}
		%Adapt labelwidth (sorry for the ugly hack)
		\setlist[description]{leftmargin=!, labelwidth=4em, before=\doublespacing}
		\IfLanguageName {finnish} {
			\printacronyms[title=Lyhenteet, nogroupskip]
		}{
			\printacronyms[title=List of Abbreviations, nogroupskip]
		}
		\setlist[description]{leftmargin=!, labelwidth=7em, before=\doublespacing}
		\vspace{1cm}
		\printglossary[nogroupskip]
		\setlist[description]{style=standard} % reset settings back to default
	}
	\addtocontents{toc}{\cftpagenumberson{chapter}}
\end{singlespacing}

\clearpage


%----------------------------------------------------------------------------------------
%	CONTENT
%----------------------------------------------------------------------------------------

\input{style/content.tex}%reset page number to 1, etc.

% Thesis content if you strictly follow the "Final Year Project guide". Of course, you
% can adapt to your specific needs (add more chapter, rename them, etc.).
% Introduction

\chapter{Johdanto}

TODO

...

Insinöörityön koodiesimerkeissä käytetään \glsdisp{ts}{TypeScriptiä} ja \glsdisp{js}{JavaScriptiä}. Vaikka kaikki validi JavaScript koodi on validia TypeScript koodia, niin kielten nimiä käytetään työssä tilanteeseensopien. Kirjoitetaan TypeScriptistä, kun sen tarjoamat tyypit JavaScriptin päälle ovat tilanteeseen nähden merkittäviä. Toisaalta kirjoitetaan JavaScriptistä, kun tyypit eivät ole esimerkille merkittäviä.

...
% uncomment what you need.
\input{chapters/projectSpec.tex}
% Material and Methods
%\clearpage%if the chapter heading starts close to bottom of the page, force a line break and remove the vertical vspace
\vspace{21.5pt}
\chapter{Toteutus}

TODO

% Theoretical background
%\clearpage % Uncomment if needed to force a page break before the chapter
\vspace{21.5pt}
\chapter{Teoreettinen tausta}

Jotta voidaan tutkia sitä, miten funktionaalinen ohjelmointi vaikuttaa ohjelmointiin \gls{ts} ympäristössä, on kartoitettava mitä ylipäätään funktionaalinen ohjelmointi on kieliagnostisesti. Tämän jälkeen voidaan tarkastella miten funktionaalinen ohjelmointi on näkynyt \gls{ts} projekteissa, ja lopuksi voidaan vertailla mikä osa tästä on kategorisesti funktionaalista, mikä olio-ohjelmallista, tai pelkästään proseduaalista.

Näillä tiedoilla voidaan lähteä tutkimaan miten käytänteet näkyvät hyödyllisesti tai haitallisesti ohjelmistoprojekteissa, joissa halutaan käyttää funktionaalista ohjelmointia pragmaattisesti.

\section{Funktionaalinen ohjelmointi}

TODO

Nopea historia, lambdakalkyyli ja kombinaattorit (Kuva \ref{code:javascript_combinators}). Olennaiset asiat. Immutability ja "pure function"

\begin{code}
  \begin{minted}{javascript}
const I  = x => x
const K  = x => y => x
const A  = f => x => f (x)
const T  = x => f => f (x)
const W  = f => x => f (x) (x)
const C  = f => y => x => f (x) (y)
const B  = f => g => x => f (g (x))
const S  = f => g => x => f (x) (g (x))
const S_ = f => g => x => f (g (x)) (x)
const S2 = f => g => h => x => f (g (x)) (h (x))
\end{minted}
  \caption{Yleiset kombinaattorit esitettynä JavaScriptissä \cite{javascript_combinators}}
  \label{code:javascript_combinators}
\end{code}

TODO

\subsection{Yhdistefunktiot ja niiden merkitys}

Funktionaalinen ohjelmointi on pääasiassa vain ja ainoastaan funktioiden yhdistämistä tai koostamista.


TODO

\begin{code}
  \begin{equation}
    f(x) = 2x
  \end{equation}
  \begin{equation}
    g(x) = x + 3
  \end{equation}
  \begin{equation}
    h(x) = f(g(x)) = 2(x + 3) = 2x + 6
  \end{equation}
  \caption{Matemaattinen esimerkki funktiokompositiosta}
  \label{equation:composition}
\end{code}
\bigskip

\begin{code}
  \begin{minted}{haskell}
f :: Int -> Int
f x = 2 * x

g :: Int -> Int
g x = x + 3

h :: Int -> Int
h x = f (g x)
\end{minted}
  \caption{Haskell-esimerkki funktiokompositiosta}
  \label{code:haskell_composition}
\end{code}
\bigskip
\begin{code}
  \begin{minted}{javascript}
function f(x) {
  return 2 * x;
}

function g(x) {
  return x + 3;
}

function h(x) {
  return f(g(x));
}
\end{minted}
  \caption{JavaScript-esimerkki funktiokompositiosta}
  \label{code:javascript_composition}
\end{code}


\subsection{Joukko-oppi}

TODO

Pari idistä siitä, että joukko-oppia pitäisi ajatella aina. Ei vain funktionaalista ohjelmointia harjoittaessa.
\begin{figure}[htbp]
  \centering
  \begin{tikzpicture}
    % Ellipses for sets A, B, and C with background colors
    \draw[thick, ellipse, fill=green!20] (0,0.5) ellipse (0.8cm and 2cm);  % Set A with light blue background
    \draw[thick, ellipse, fill=blue!20] (3,0.5) ellipse (0.8cm and 2cm); % Set B with light green background
    \draw[thick, ellipse, fill=red!20] (6,0.5) ellipse (0.8cm and 2cm);   % Set C with light red background

    % Nodes for set A (without borders)
    \node[draw=none] (A1) at (0,2) {1};
    \node[draw=none] (A2) at (0,1) {2};
    \node[draw=none] (A3) at (0,0) {3};
    \node[draw=none] (A4) at (0,-1) {4};

    % Nodes for set B (without borders)
    \node[draw=none] (B1) at (3,2) {1};
    \node[draw=none] (B2) at (3,1) {2};
    \node[draw=none] (B3) at (3,0) {3};
    \node[draw=none] (B4) at (3,-1) {4};

    % Nodes for set C (without borders)
    \node[draw=none] (C1) at (6,2) {1};
    \node[draw=none] (C2) at (6,1) {2};
    \node[draw=none] (C3) at (6,0) {3};
    \node[draw=none] (C4) at (6,-1) {4};

    % Arrows from A to B
    \draw[->] (A1) -- (B1);
    \draw[->] (A2) -- (B3);
    \draw[->] (A3) -- (B2);
    \draw[->] (A4) -- (B4);

    % Arrows from B to C
    \draw[->] (B1) -- (C2);
    \draw[->] (B2) -- (C1);
    \draw[->] (B3) -- (C3);
    \draw[->] (B4) -- (C4);

    % Labels for sets
    \node[draw=none, text=green!50!black] at (0,3) {$A = \{1,2,3,4\}$};
    \node[draw=none, text=blue] at (3,3) {$B = \{1,2,3,4\}$};
    \node[draw=none, text=red] at (6,3) {$C = \{1,2,3,4\}$};

    % Arrow to new diagram
    \draw[->, line width=1.5mm, black] (3,-1.75) -- (3,-2.75) node[midway, right] {};

    % New diagram with direct A to C composition
    % Ellipses for sets A and C below the original picture
    \draw[thick, ellipse, fill=green!20] (0,-4.5) ellipse (0.8cm and 2cm);  % Set A with light blue background
    \draw[thick, ellipse, fill=red!20] (6,-4.5) ellipse (0.8cm and 2cm);   % Set C with light red background

    % Nodes for set A (without borders) - Lower position
    \node[draw=none] (A1b) at (0,-3) {1};
    \node[draw=none] (A2b) at (0,-4) {2};
    \node[draw=none] (A3b) at (0,-5) {3};
    \node[draw=none] (A4b) at (0,-6) {4};

    % Nodes for set C (without borders) - Lower position
    \node[draw=none] (C1b) at (6,-3) {1};
    \node[draw=none] (C2b) at (6,-4) {2};
    \node[draw=none] (C3b) at (6,-5) {3};
    \node[draw=none] (C4b) at (6,-6) {4};

    % Direct arrows from A to C (composition)
    \draw[->] (A1b) -- (C2b);
    \draw[->] (A2b) -- (C3b);
    \draw[->] (A3b) -- (C1b);
    \draw[->] (A4b) -- (C4b);

  \end{tikzpicture}
  \caption{Funktioiden yhdistäminen joukko-opin näkökulmasta. Funktiot f : A$\,\to\,$B ja g : B$\,\to\,$C yhdistämällä saadaan g : A$\,\to\,$C.}
\end{figure}


Kelaa mitkä chancet tehä juttuja tehokkaaksi fiksulla kääntäjällä! Voidaan poistaa turhaa laskentaa.

\subsection{Kategoriateoria}

TODO

Pari idistä siitä, mitä hyötyy kun lähtee syvään päähänn. Toisaalta myös se, että miten sinne edes pääsee ja miten sen voi erottaa.

\section{Tyypitys}

TODO

Tyypitys on todella tärkeää. Tyypit ohjaavat koodia ja estävät vääriä tiloja ohjemakoodin kääntämisvaiheessa.

TODO Mites \gls{js} :D



\subsection{Parametrinen polymorfismi}

TODO

Haskelille ominaista, \gls{fp} ominaista.


\subsection{Hindley-Milner tyyppijärjestelmä}

TODO

Feldmanin mukaan, mahdottomat tilat ovat testausta parempaa \cite{impossiblebetter}.

Myös Dijkstran idis abstraktiosta sopii kauniisti tähän, ja esimerkki siitä et mitä \gls{fp} on parhaimmillaan.

Parametrinen polymorfismi näkyy täs myös.

\section{Pragmaattisuus}

TODO

Kun funktionaalisen ohjelmoinnin kauneus on matemaattisissa kulmakivissä jotka eivät taivu, miten sitä voi yhdistää nykymaailman pikakoodaukseen? Python, \gls{ts} yms. Chat gpt? Miten se edes osaa kirjoittaa fp koodia :D

TODO miten uncle bob sanoo? \cite{martin2017pragmaticfp}

\subsection{A monad is a monoid in the category of endofunctors}

TODO

Hyvää bäntteriä siitä, miksei fp ole lähtenyt lentoon, miksi pragmaattisuus on hot. Miten fp näkyy nykymaailmassa ja miten se mielletään. Pohja kategoriateoriassa pelottaa.

Nopea kattaus \textquote{developer advocate} -heppujen mielipiteistä.

Viittaus Promise-spesifikaation monadittomuuteen, mainitaan myös alempana

\subsection{Funktionaalisen ohjelmoinnin ripottelu}

TODO

Idiksiä siitä että juuri otetaan funktionaalisen ohjemoinnin parhaat palat

Youtube videoista konkreettisia esimerkkejä.

\section{Funktionaalinen ohjelmointi \& Typescript}


TODO

Typescriptiä tullaan käyttämään insinöörityön koodiesimerkeissä, niin on aiheellista näyttää miten Typescriptissä \gls{fp} näkyy.

TODO TypeScript tyypit HYVÄ! Rich Hickey tykkää myös, voi tehä "product type" tms

Toisaalta TS -> overhead, välillä joutuu kikkailemaan

\subsection{Promise-tietorakenne ja osittaismonadisuus}

Promise ei ole monadi :( \cite{read-it-later-11481,promises-spec-94}

Se ei ole monadi, koska \gls{js} spesifikaatiota määriteltäessä ei päästy yhteisymmärrykseen siitä, että Promise-tietorakenteen tulisi olla monadi. Syy on moniuloitteinen, mutta pääasiallisesti sivuaa sitä, että funktionaalinen ohjelmointi ja sen käsitteet ovat vaikeita selittää kohtuullisella tarkkuudella. \citep{promises-spec-94}

Tämä kirjoitushetkellä jo yli kymmenen vuotta sitten tehty päätös heikensi merkittävästi \gls{js}:n mahdollisuuksia viedä funktionaalista ohjelmointia eteenpäin. \Gls{js} kielenä on laajuudeltaan ongelmallinen. Uusia ominaisuuksia ei voi lisätä, jos vastaavanlaisia ominaisuuksia on aiemmin ollut osana suosittuja \gls{js}-kirjastoja \cite{proposal-joint-iteration,prototype_library_trends}.

TODO lisää siitä että ei tulla ikinä saamaan "builtin" monadeja
TODO myös että no mitä se tarkoittaa, onko juuri promisen kannalta mitään merkitystä

\subsection{Ramda.js}

Koska funktionaalisen ohjelmakoodin tuottaminen \gls{js} tai \gls{ts} ympäristössä enemmän tai vähemmän vaatii jotain kirjastoa, joko itse ohjelmoitua tai ulkoista, niin tähän insinöörityöhön on kirjastoksi valittu Ramda.js.

Ramda.js, ja kirjaston taustalla oleelinen Fantasy Land -algebrallinen spesifikaatio, ovat kirjastoja jotka pyrkivät johdonmukaiseen funktionaaliseen ohjelmointiin JavaScriptissä. Näiden kirjastojen kehittäjistöstä löytyy samoja henkilöitä jotka aikoinaan pyrkivät saamaan Promise-tietorakenteen virallisesti monadiseksi \cite{ramda:contributors,fantasy-land:contributors,promises-spec-94}. Näyttää siltä, että kun tähän tavoitteeseen ei päästy, päättivät kehittäjät tuoda funktionaalisen ohjelmoinnin JavaScriptiin omin keinoin.

TODO miten ramda toimii

TODO miten ramda sopii omien funktioiden kanssa

TODO miten typescript? ei hyvä, tyypit ovat mitä sattuu (kokemuspohjaisesti)

\subsection{Rajoitteet}

TODO

Mitkä asiat eivät ole mahdollisia, tai mitkä ovat erittäin hankalia / järjettömiä kun ohjelmoidaan \gls{fp} koodia \gls{ts}-maailmassa.


\section{Vertailu yhteenvetona}

TODO

Ehkä nopea yhteenveto että mikä on täysin \gls{fp} ominaista ja mikä ei. Mahdollisesti jokin kiva taulukko.
% Material and Methods
%\clearpage%if the chapter heading starts close to bottom of the page, force a line break and remove the vertical vspace
\vspace{21.5pt}
\chapter{Toteutus}

TODO

\section{Hyödyt ja haitat}
TODO

% Material and Methods
%\clearpage%if the chapter heading starts close to bottom of the page, force a line break and remove the vertical vspace
\vspace{21.5pt}
\chapter{Johtopäätös}

TODO

\section{Pohdinta}

TODO
% fp hyvä, kieliagnostiset datatyypit ja rakenteet hyvä

% Kikkailu huono.

% fp koodin kirjoittaminen vain koska "fp hyvä" on huono

\section{Suositukset}

TODO

\section{Henkilökohtaiset reflektiot}

TODO

% Miten onnistua kirjoittamaan sitä kuuluisaa funktionaalista koodia sitten loppujen lopuksi?

% Varmaan siten, että yrittää miettiä kategoriateoriaa. Kategoriateoria on kovin teoreettista, mutta teoreettinen = hyvä pohja, jos teoria on oikeasti totta eikä harhaunta.

\section{Loppusanat}

%----------------------------------------------------------------------------------------
%	BIBLIOGRAPHY REFERENCES
%----------------------------------------------------------------------------------------

% Bibliography.
% Normally you do not edit this file.
% To add bibliography in your text, add them first in the biblio.bib file and
% reference them with the \cite{} command in your text. Check also \citeauthor
% and \citeyear (e.g. for copied figure caption), the \cites for multi-cites
% with page numbers,...

\clearpage
%line space
%\singlespacing %removed otherwise the appendix are also single space
\begin{singlespacing}
    \vspace*{-16pt}
    \printbibliography[title=\IfLanguageName{finnish}{Lähteet}{References}]
\end{singlespacing}

%for conting the pages
\label{LastPage}~


%----------------------------------------------------------------------------------------
%	APPENDICES
%----------------------------------------------------------------------------------------

% % Appendix
% Normally, you do not have to edit this file.

%start appendix
\appendix
%no page number for appendix in table of content
\addtocontents{toc}{\cftpagenumbersoff{chapter}}
%appendix sections and subsections not in table of content
\settocdepth{chapter}
%add "Appendices" in the table of content
\addappheadtotoc
%have Appendix 1 (instead of Appendix A)
\renewcommand{\thechapter}{\arabic{chapter}}

\newcommand\liite[1]{
    %each appendix start on a new page
    \clearpage
    %each appendix restart page num to one
    \setcounter{page}{1}
    %special counter for appendix TODO: this is a ugly quick hack :( Should find a better way to count the page per appendix.
    \newtotcounter{appx#1}
    %overwrite the header
    \makeevenhead{plain}{}{}{\appname \thechapter \\ \thepage\,(\stepcounter{appx#1}\total{appx#1})}
    \makeoddhead{plain}{}{}{\appname \thechapter \\ \thepage\,(\stepcounter{appx#1}\total{appx#1})}}



% %force smaller vertical spacing in table of content
% %!!! There can be some fun depending if the appendices have (sub)sections or not :D
% % You will have to play with these numbers and eventually add the \vspace line  before
% % some \chapter and force another number.
% % To add more fun, time to time the table of content get wrong after a build :(
% \addtocontents{toc}{\vspace{11pt}}
% \pretocmd{\chapter}{\addtocontents{toc}{\protect\vspace{-24pt}}}{}{}

% \liite{1}% This is a hack to have right page numbering for each appendix. Make sure to
% % use a unique number for each appendix.
% \input{sample/Xappendix1.tex}% Sample content to demonstrate appendix in LaTeX. You
% % are safe to delete this lines (and the next samples) once you refreshed your LaTeX
% % skills (and safe to delete the sample folder and all its file too).

% %\addtocontents{toc}{\vspace{11pt}}%fix vertical space for Table of Content
% \liite{2}
% \input{sample/Xappendix2.tex}

% \addtocontents{toc}{\vspace{11pt}}
% \liite{3}
% \input{sample/X_R_example.tex}


%----------------------------------------------------------------------------------------
%	THIS IS THE END
%----------------------------------------------------------------------------------------
\end{document}
